% !TEX program = xelatex
\documentclass[11pt,a4paper]{article}

%% ── Язык и шрифты ────────────────────────────────────────────────────────────
\usepackage{fontspec}
\usepackage{polyglossia}
\setmainlanguage{russian}
\setotherlanguage{english}
\setmainfont{Times New Roman}[
  BoldFont      = Liberation Serif Bold,
  ItalicFont    = Liberation Serif Italic,
  BoldItalicFont= Liberation Serif Bold Italic
]
\setsansfont{PT Sans}
\setmonofont{PT Mono}

%% ── Страница ─────────────────────────────────────────────────────────────────
\usepackage[top=2.5cm, bottom=2.5cm, left=2.5cm, right=2cm]{geometry}
\widowpenalty=10000
\clubpenalty=10000

%% ── Графика и цвета ──────────────────────────────────────────────────────────
\usepackage{graphicx}
\usepackage[table]{xcolor}
\usepackage{colortbl}
\definecolor{mainblue}{RGB}{41, 128, 185}
\definecolor{darkblue}{RGB}{31, 97, 141}
\definecolor{lightgray}{RGB}{245, 245, 245}
\definecolor{midgray}{RGB}{189, 195, 199}
\definecolor{accent}{RGB}{231, 76, 60}
\definecolor{green}{RGB}{39, 174, 96}

%% ── Таблицы ──────────────────────────────────────────────────────────────────
\usepackage{booktabs}
\usepackage{tabularx}
\usepackage{multirow}
\usepackage{array}
\newcolumntype{C}[1]{>{\centering\arraybackslash}p{#1}}
\newcolumntype{L}[1]{>{\raggedright\arraybackslash}p{#1}}
\newcolumntype{R}[1]{>{\raggedleft\arraybackslash}p{#1}}

%% ── Окружения и оформление ────────────────────────────────────────────────────
\usepackage{float}
\usepackage{caption}
\usepackage[most]{tcolorbox}
\usepackage{enumitem}
\setlist[itemize]{nosep, leftmargin=*, topsep=3pt, partopsep=0pt}
\setlist[enumerate]{nosep, leftmargin=*, topsep=3pt, partopsep=0pt}

%% ── Заголовки разделов ────────────────────────────────────────────────────────
\usepackage{titlesec}
\newcommand{\sectionrule}{\vspace{-2pt}{\color{mainblue}\titlerule[0.6pt]}\vspace{4pt}}
\titleformat{\section}{\large\bfseries\color{darkblue}}{\thesection.}{0.5em}{}[\sectionrule]
\titleformat{\subsection}{\normalsize\bfseries\color{darkblue}}{\thesubsection.}{0.5em}{}
\titleformat{\subsubsection}{\normalsize\bfseries}{\thesubsubsection.}{0.5em}{}
\titlespacing*{\section}{0pt}{14pt}{6pt}
\titlespacing*{\subsection}{0pt}{10pt}{4pt}
\titlespacing*{\subsubsection}{0pt}{8pt}{3pt}

%% ── Колонтитулы ─────────────────────────────────────────────────────────────
\usepackage{fancyhdr}
\setlength{\headheight}{14pt}
\pagestyle{fancy}
\fancyhf{}
\fancyhead[L]{\small\color{darkblue}Несупервизируемая сегментация ДЗЗ}
\fancyhead[R]{\small\color{darkblue}DINOv2 + Mean Teacher}
\fancyfoot[C]{\small\thepage}
\renewcommand{\headrulewidth}{0.4pt}
\renewcommand{\headrule}{\color{mainblue}\hrule width\headwidth height\headrulewidth}

%% ── Математика ───────────────────────────────────────────────────────────────
\usepackage{amsmath, amssymb}

%% ── Гиперссылки ──────────────────────────────────────────────────────────────
\usepackage[colorlinks=true, linkcolor=darkblue, urlcolor=mainblue, citecolor=accent]{hyperref}

%% ── Библиография ─────────────────────────────────────────────────────────────
\usepackage[numbers, sort&compress]{natbib}

%% ── Метаданные ───────────────────────────────────────────────────────────────
\hypersetup{
  pdftitle={Несупервизируемая семантическая сегментация снимков ДЗЗ},
  pdfauthor={DanilaE},
  pdfsubject={Deep Learning, Remote Sensing, Semantic Segmentation}
}

%% ── Пути к рисункам ─────────────────────────────────────────────────────────
\graphicspath{{/data/NNDP-Lab/report_figures/}}

%% ── tcolorbox стили ──────────────────────────────────────────────────────────
\tcbuselibrary{breakable, skins}
\newtcolorbox{infobox}[1][]{
  colback=mainblue!6, colframe=mainblue!70,
  fonttitle=\bfseries\small, title=#1,
  left=6pt, right=6pt, top=4pt, bottom=4pt,
  sharp corners=south
}
\newtcolorbox{resultbox}[1][]{
  colback=green!6, colframe=green!70,
  fonttitle=\bfseries\small, title=#1,
  left=6pt, right=6pt, top=4pt, bottom=4pt,
  sharp corners=south
}
\newtcolorbox{warnbox}[1][]{
  colback=accent!5, colframe=accent!60,
  fonttitle=\bfseries\small, title=#1,
  left=6pt, right=6pt, top=4pt, bottom=4pt,
  sharp corners=south
}

%%=============================================================================
\begin{document}
\thispagestyle{empty}

%% ── Титульная страница ───────────────────────────────────────────────────────
\begin{center}
  \vspace*{1.5cm}
  {\color{mainblue}\rule{\textwidth}{2pt}}\\[0.6cm]
  {\Huge\bfseries\color{darkblue}
    Несупервизируемая семантическая\\[6pt]
    сегментация снимков\\[6pt]
    дистанционного зондирования Земли}\\[0.8cm]
  {\Large\color{mainblue}
    Трансферное обучение с псевдометками\\
    и самообучением (Mean Teacher)}\\[0.6cm]
  {\color{mainblue}\rule{\textwidth}{2pt}}\\[1.5cm]

  \begin{tabular}{rl}
    \textbf{Автор:}   & DanilaE \\[4pt]
    \textbf{Датасет:} & LoveDA (Kaggle / torchgeo) \\[4pt]
    \textbf{Дата:}    & 26 мая 2026 г. \\[4pt]
    \textbf{Устройство:} & NVIDIA GeForce RTX~5070~Ti (16\,ГБ VRAM) \\[4pt]
    \textbf{Фреймворк:} & PyTorch 2.12.1+cu132, DINOv2 ViT-B/14 \\
  \end{tabular}

  \vspace{1.8cm}

  \begin{minipage}{0.85\textwidth}
    \begin{tcolorbox}[colback=darkblue!8, colframe=darkblue, title=Аннотация, fonttitle=\bfseries]
      Представлен полный пайплайн несупервизированной семантической сегментации спутниковых
      снимков без единой размеченной выборки. Замороженный DINOv2 ViT-B/14 инициализирует
      псевдометки через MiniBatch K-means (2\,582\,528 патч-векторов, $d{=}768$).
      Декодер SegFormer обучается итеративно через механизм Mean Teacher:
      teacher генерирует псевдометки при слабой аугментации, student обучается на
      сильных аугментациях. Эксперименты проведены на датасете LoveDA (2\,522 обучающих
      + 1\,669 валидационных снимков, 7 классов земного покрова). За полное обучение
      (55 эпох, 2\,ч\,18\,мин) потери снизились с 1.60 до 0.15 ($-91\,\%$);
      лучший val mIoU составил \textbf{21.38\,\%}.
    \end{tcolorbox}
  \end{minipage}
\end{center}

\vfill
\begin{center}
  \small\color{midgray}NNDP-Lab $\cdot$ 2026
\end{center}

\newpage

%% ── Оглавление ───────────────────────────────────────────────────────────────
\tableofcontents
\newpage

%%=============================================================================
\section{Введение и постановка задачи}

Семантическая сегментация снимков дистанционного зондирования Земли (ДЗЗ) является
одной из ключевых задач геоинформатики. Её применяют для мониторинга землепользования,
планирования городской инфраструктуры, оценки ущерба от стихийных бедствий и точного
земледелия. Традиционные методы сегментации требуют большого объёма вручную размеченных
данных, получение которых для спутниковых снимков трудозатратно и дорогостояще:
эксперты-интерпретаторы должны работать с высокоразрешающими мультиспектральными
изображениями, охватывающими площади в сотни квадратных километров.

\textbf{Цель работы} — разработать систему семантической сегментации снимков ДЗЗ,
не требующую ни одной размеченной выборки на этапе обучения. Для этого используется
комбинация трёх взаимодополняющих техник. \textbf{Трансферное обучение} на основе
предобученного визуального трансформера DINOv2~\cite{oquab2024dinov2} (Meta AI, 2023),
обученного методом самообучения на 142~млн изображений, обеспечивает богатые
семантические признаки без дополнительной разметки. \textbf{Кластеризация признаков}
через MiniBatch K-means инициализирует начальные псевдометки непосредственно из
DINOv2-представлений. \textbf{Self-training с Mean Teacher}~\cite{tarvainen2017mean}
итеративно уточняет псевдометки, обучая student-модель согласовываться с teacher-моделью,
веса которой обновляются как экспоненциальное скользящее среднее (EMA) весов student.

\begin{infobox}[Датасет]
  \textbf{LoveDA}~\cite{wang2021loveda} (Land-cOVEr Domain Adaptive semantic segmentation)
  распространяется через Kaggle и библиотеку torchgeo. Содержит высокоразрешающие
  (0.3\,м/пиксель) RGB-снимки городской и сельской местности Китая с пиксельной
  разметкой по 7 классам. Общий объём обучающей и валидационной выборок — 4\,191
  изображение размером $1024 \times 1024$ пикселей.
\end{infobox}

%%=============================================================================
\section{Данные}

\subsection{Датасет LoveDA}

LoveDA охватывает два контрастных домена --- городской (Urban) и сельский (Rural) ---
что делает его актуальным для исследования переноса домена. Все снимки сделаны с
платформы Google Earth и имеют единое пространственное разрешение 0.3\,м/пиксель.

\begin{table}[htbp]
\centering
\caption{Основные характеристики датасета LoveDA}
\label{tab:loveda_stats}
\small
\begin{tabularx}{0.88\textwidth}{L{4.8cm} X}
  \toprule
  \textbf{Параметр} & \textbf{Значение} \\
  \midrule
  Источник & Google Earth (Китай: Нанкин, Чэнду, Ухань) \\
  Пространственное разрешение & 0.3\,м/пиксель \\
  Спектральные каналы & RGB (3 канала) \\
  Размер снимка & $1024 \times 1024$ пикселей \\
  Обучающих снимков & 2\,522 (Urban: 1\,156, Rural: 1\,366) \\
  Валидационных снимков & 1\,669 (Urban: 677, Rural: 992) \\
  Тестовых снимков & 1\,796 (разметка недоступна) \\
  Число классов & 7 \\
  \bottomrule
\end{tabularx}
\end{table}

\begin{figure}[htbp]
  \centering
  \includegraphics[width=0.96\textwidth]{fig1_samples.png}
  \caption{Примеры изображений из валидационной выборки LoveDA (Urban).
           Слева направо: исходный снимок, маска Ground Truth, наложение классов,
           гистограмма пикселей по классам.}
  \label{fig:samples}
\end{figure}

На рис.~\ref{fig:samples} представлены репрезентативные примеры снимков и соответствующих
масок GT. Каждый цвет соответствует одному из семи классов земного покрова: фон
(серый), здания (красно-оранжевый), дороги (жёлтый), вода (синий), пустошь
(коричневый), лес (зелёный), сельскохозяйственные угодья (золотой).

\subsection{Распределение классов}

Подсчёт проводился по всем 4\,191 снимкам обучающей и валидационной выборок:

\begin{table}[htbp]
\centering
\caption{Реальное распределение пикселей по классам (Train+Val, $\approx 4.4$ млрд пикселей)}
\label{tab:class_dist}
\small
\rowcolors{2}{lightgray}{white}
\begin{tabularx}{0.90\textwidth}{l C{1.8cm} C{1.8cm} C{1.8cm}}
  \toprule
  \textbf{Класс} & \textbf{Urban} & \textbf{Rural} & \textbf{Итого} \\
  \midrule
  Background (Фон)    & 38.3\,\% & 31.9\,\% & \textbf{34.7\,\%} \\
  Building (Здания)   & 17.2\,\% &  3.0\,\% & 9.2\,\% \\
  Road (Дороги)       &  8.3\,\% &  2.2\,\% & 4.9\,\% \\
  Water (Вода)        &  6.4\,\% &  9.7\,\% & 8.2\,\% \\
  Barren (Пустошь)    &  6.5\,\% &  3.3\,\% & 4.7\,\% \\
  Forest (Лес)        &  8.6\,\% & 14.9\,\% & \textbf{12.2\,\%} \\
  Agriculture (С/х)   & 14.7\,\% & 34.9\,\% & \textbf{26.1\,\%} \\
  \bottomrule
\end{tabularx}
\end{table}

Класс Background доминирует (34.7\,\%), Agriculture занимает второе место (26.1\,\%).
Датасет существенно несбалансирован: Barren встречается в $7\times$ реже Background.
Резкий контраст между доменами --- Urban отличается высокой долей зданий и дорог,
Rural --- леса и сельскохозяйственных угодий --- подтверждает необходимость совместного
обучения на обоих доменах.

\begin{figure}[htbp]
  \centering
  \includegraphics[width=0.96\textwidth]{fig3_urban_rural.png}
  \caption{Визуальное сравнение Urban и Rural доменов: исходные изображения и маски GT.
           Городские снимки характеризуются высокой плотностью зданий и дорог,
           сельские --- обширными сельскохозяйственными угодьями и лесами.}
  \label{fig:urban_rural}
\end{figure}

\subsection{Предобработка}

Изображения масштабируются до фиксированного размера (224 или 448\,пикселей по стороне)
и нормализуются по статистике ImageNet: $\mu = (0.485, 0.456, 0.406)$,
$\sigma = (0.229, 0.224, 0.225)$. Размер выбирается кратным 14 --- размеру патча
DINOv2 ViT, что обеспечивает целочисленное разрешение карт признаков.
Разметка при обучении \textbf{не используется}; маски GT применяются исключительно
для финальной оценки через венгерский алгоритм.

%%=============================================================================
\section{Методология}

\subsection{Архитектура системы}

\begin{figure}[htbp]
  \centering
  \includegraphics[width=0.88\textwidth]{fig5_pipeline.png}
  \caption{Пайплайн несупервизируемой сегментации. Слева направо: обработка входного
           изображения DINOv2 backbone, декодирование SegFormer, трёхфазное обучение
           (K-means $\to$ warm-up $\to$ self-training с Mean Teacher).}
  \label{fig:pipeline}
\end{figure}

Система состоит из двух компонентов, объединённых в пару student/teacher:

\textbf{DINOv2 ViT Backbone.} Принимает RGB-изображение размером $H \times W$ и
извлекает признаки четырёх промежуточных слоёв трансформера, каждый с пространственным
разрешением $H/14 \times W/14 \times D$ (для ViT-S/14: $D = 384$). При замороженном
backbone обучается только декодер; с раунда 3 размораживаются последние четыре
блока трансформера.

\textbf{SegFormer MLP Decoder~\cite{xie2021segformer}.} Четыре независимых MLP
проецируют признаки каждого масштаба в пространство размерности \texttt{decoder\_dim}
(по умолчанию 256). Все карты признаков билинейно интерполируются до размера наибольшей
карты, конкатенируются и сливаются свёрткой $1{\times}1$ + BN + ReLU. Финальная
свёртка $1{\times}1$ выдаёт тензор логитов формы \texttt{num\_classes} $\times H \times W$.

\subsection{Трёхфазный пайплайн обучения}

\subsubsection{Фаза 0 --- Инициализация псевдометок (K-means)}

Замороженный backbone извлекает патч-признаки последнего слоя для всех обучающих снимков.
Все $N_\text{img} \times H_p \times W_p$ векторов (где $H_p = W_p = H/14$) подаются в
MiniBatch K-means с $k{=}7$ кластерами. Каждому патчу присваивается кластерный ID,
образуя начальную псевдометку с разрешением $H_p \times W_p$.

\subsubsection{Фаза 1 --- Разогрев (Warm-up)}

Декодер обучается на K-means псевдометках при замороженном backbone стандартной
кросс-энтропией по всем пикселям. Цель фазы --- привести декодер в разумное начальное
состояние перед итеративным самообучением.

\subsubsection{Фаза 2 --- Self-training раунды}

Каждый раунд включает четыре этапа. \textbf{Обновление псевдометок:} teacher-модель
прогоняет все обучающие снимки и сохраняет предсказания для пикселей, где максимальная
softmax-вероятность превышает порог $\theta$. \textbf{Обучение student:} teacher
получает слабо аугментированные снимки (случайное горизонтальное отражение) и формирует
мягкие псевдометки с маской уверенности; student получает сильно аугментированные снимки
(случайный crop + colour jitter + Gaussian blur). Суммарная функция потерь:

\begin{equation}
  \mathcal{L} = \mathcal{L}_\text{CE}(s_\theta, \hat{y}, \mathbf{1}) +
                0.5 \cdot \mathcal{L}_\text{KL}(s_\theta, t_\theta, \mathbf{m}_\text{conf}),
  \label{eq:loss}
\end{equation}

\noindent где $\mathcal{L}_\text{CE}$ --- кросс-энтропия по всем пикселям,
$\mathcal{L}_\text{KL}$ --- дивергенция Кульбака--Лейблера только по уверенным пикселям
($\mathbf{m}_\text{conf}$). \textbf{EMA-обновление teacher:} после каждого шага

\begin{equation}
  \theta_T \leftarrow \alpha \cdot \theta_T + (1 - \alpha) \cdot \theta_S, \quad \alpha = 0.999.
  \label{eq:ema}
\end{equation}

\noindent \textbf{Снижение порога:} каждый раунд $\theta$ уменьшается на $\Delta\theta = 0.025$,
но не ниже $\theta_\text{min} = 0.70$, постепенно включая менее уверенные предсказания.

\subsection{Оценка качества}

Поскольку кластерные ID произвольны и не совпадают с семантическими классами,
для оценки применяется венгерский алгоритм (Hungarian matching / linear assignment),
находящий оптимальную перестановку кластеров:

\begin{equation}
  \text{mIoU} = \frac{1}{C} \sum_{c=1}^{C} \text{IoU}\!\left(\sigma^*(c), c\right),
  \label{eq:miou}
\end{equation}

\noindent где $\sigma^*$ --- оптимальная перестановка, $C = 7$. Это единственный
корректный способ измерить качество несупервизированной сегментации.

%%=============================================================================
\section{Эксперименты и результаты}

\subsection{Конфигурации и гиперпараметры}

\begin{table}[htbp]
\centering
\caption{Сравнение конфигураций обучения}
\label{tab:configs}
\small
\begin{tabularx}{0.92\textwidth}{L{2.8cm} C{2cm} C{1.5cm} C{1.8cm} C{2.5cm}}
  \toprule
  \textbf{Режим} & \textbf{Backbone} & \textbf{Разреш.} & \textbf{Batch} & \textbf{Устройство} \\
  \midrule
  Быстрый тест                      & ViT-S/14          & 224\,px          & 16        & RTX 5070\,Ti \\
  Полный CPU                        & ViT-S/14          & 224\,px          &  8        & CPU \\
  \textbf{Полный GPU (этот запуск)} & \textbf{ViT-B/14} & \textbf{448\,px} & \textbf{16} & \textbf{RTX 5070\,Ti} \\
  \bottomrule
\end{tabularx}
\end{table}

Гиперпараметры полного GPU-запуска: оптимизатор AdamW, LR декодера $10^{-4}$,
LR backbone $10^{-5}$ (разморозка с раунда~3), weight decay $10^{-4}$,
EMA decay $\alpha = 0.999$, 5 эпох warm-up, 5 раундов self-training по 10 эпох,
начальный порог уверенности $\theta_0 = 0.95$ с убыванием 0.025 за раунд (минимум 0.70).

\begin{table}[htbp]
\centering
\caption{Сравнение двух экспериментов (быстрый тест vs. полное обучение)}
\label{tab:compare}
\small
\begin{tabularx}{0.96\textwidth}{L{3.8cm} C{2.2cm} C{2.2cm} C{2.0cm} C{2.0cm}}
  \toprule
  \textbf{Параметр} & \textbf{Быстрый тест} & \textbf{Полное GPU} \\
  \midrule
  Backbone        & ViT-S/14   & ViT-B/14 \\
  Разрешение      & 224\,px    & 448\,px \\
  Признаки K-means & 645\,632 ($d{=}384$) & 2\,582\,528 ($d{=}768$) \\
  Warm-up         & 1 эпоха    & 5 эпох \\
  Self-training   & 2 раунда $\times$ 2 эп. & 5 раундов $\times$ 10 эп. \\
  Порог $\theta$  & 0.70 (фикс.) & 0.95 $\to$ 0.85 \\
  Время обучения  & $\approx$10\,мин & 2\,ч\,18\,мин \\
  Loss (нач. $\to$ кон.) & 1.60 $\to$ 0.44 & 1.60 $\to$ 0.15 \\
  \textbf{Best val mIoU} & \textbf{9.86\,\%} & \textbf{21.38\,\%} \\
  \bottomrule
\end{tabularx}
\end{table}

\subsection{Реальные результаты обучения на GPU}

\begin{figure}[htbp]
  \centering
  \includegraphics[width=0.88\textwidth]{fig6_training_curves.png}
  \caption{Динамика обучения на RTX~5070~Ti (полный запуск, ViT-B/14, 448\,px).
           Слева: потери убывают с 1.60 до 0.15 ($-91\,\%$) за 55 эпох.
           Справа: val mIoU с венгерским выравниванием;
           пунктирная линия --- случайный baseline ($1/7 \approx 14.3\,\%$).}
  \label{fig:curves}
\end{figure}

\begin{table}[htbp]
\centering
\caption{Метрики полного GPU-прогона (ViT-B/14, 448\,px, 55 эпох, RTX~5070~Ti)}
\label{tab:results}
\small
\rowcolors{2}{lightgray}{white}
\begin{tabularx}{0.96\textwidth}{L{3.4cm} C{1.5cm} C{2.2cm} C{2.8cm} C{2.5cm}}
  \toprule
  \textbf{Этап} & \textbf{Эпохи} & \textbf{Train Loss} & \textbf{Val mIoU} & \textbf{Уверен. пиксели} \\
  \midrule
  Warm-up       & 5  & 1.60 $\to$ 1.50 & \textbf{21.38\,\%} & --- \\
  Round 1 ($\theta{=}0.95$) & 10 & 1.13 $\to$ 1.02 & 14.05\,\% & 0.0\,\% \\
  Round 2 ($\theta{=}0.93$) & 10 & 0.74 $\to$ 0.70 & 9.25\,\%  & 0.0\,\% \\
  Round 3 ($\theta{=}0.90$) & 10 & 0.50 $\to$ 0.40 & 7.83\,\%  & 12.6\,\% \\
  Round 4 ($\theta{=}0.88$) & 10 & 0.28 $\to$ 0.21 & 8.44\,\%  & 59.3\,\% \\
  Round 5 ($\theta{=}0.85$) & 10 & 0.16 $\to$ 0.15 & 8.56\,\%  & 82.9\,\% \\
  \midrule
  \textbf{Best} & --- & --- & \textbf{21.38\,\%} & --- \\
  \bottomrule
\end{tabularx}
\end{table}

\begin{resultbox}[Ключевые наблюдения]
  \begin{itemize}
    \item Loss убывает с 1.60 до 0.15 ($-91\,\%$) за 55 эпох.
    \item Лучший результат достигнут \textbf{после warm-up} (21.38\,\%) ---
          до начала self-training раундов. Это вдвое превышает результат
          быстрого теста (9.86\,\%), главным образом за счёт ViT-B/14
          (dim 768 vs.\ 384) и разрешения 448\,px.
    \item В раундах 1--2 уверенных пикселей 0\,\% (порог $\theta{=}0.95$
          слишком высок), поэтому consistency loss не активен и модель
          обучается только по K-means псевдометкам, постепенно теряя
          начальное выравнивание.
    \item С раунда 3 ($\theta{=}0.90$, 12.6\,\% уверенных пикселей)
          consistency loss активируется; к раунду 5 --- 82.9\,\% уверенных
          пикселей. Loss продолжает падать, однако mIoU не восстанавливается
          до уровня warm-up --- типичное поведение при несбалансированном
          начальном пороге.
  \end{itemize}
\end{resultbox}

\subsection{Визуальное качество предсказаний}

\begin{figure}[htbp]
  \centering
  \includegraphics[width=0.88\textwidth]{fig7_predictions.png}
  \caption{Предсказания модели на валидационных снимках LoveDA (полное обучение, ViT-B/14).
           Слева: исходное изображение. По центру: Ground Truth. Справа: предсказание
           модели (разметка при обучении не использовалась).}
  \label{fig:predictions}
\end{figure}

После полного обучения модель уверенно выделяет водные объекты (синий),
лесной покров (зелёный) и здания (красно-оранжевый). Точное выравнивание
кластерных ID с семантическими классами обеспечивается венгерским алгоритмом при
финальной оценке.

\subsection{Распределение классов в датасете}

\begin{figure}[htbp]
  \centering
  \includegraphics[width=0.84\textwidth]{fig2_class_dist.png}
  \caption{Доля пикселей каждого класса по всем 4\,191 снимкам (Train+Val).
           Левая диаграмма: абсолютные доли; правая: круговая диаграмма.}
  \label{fig:dist}
\end{figure}

\begin{figure}[htbp]
  \centering
  \includegraphics[width=0.84\textwidth]{fig4_domain_dist.png}
  \caption{Распределение классов отдельно по Urban и Rural доменам.
           Контраст доменов подтверждает необходимость совместного обучения.}
  \label{fig:domain_dist}
\end{figure}

\subsection{Анализ ошибок и ограничения}

\begin{warnbox}[Ограничения и наблюдения]
  \textbf{Деградация mIoU при self-training.} Warm-up дал лучший результат (21.38\,\%).
  При $\theta_0{=}0.95$ первые два раунда работают без consistency loss (0\,\%
  уверенных пикселей), что ведёт к постепенному смещению псевдометок.
  Рекомендуется начинать с $\theta_0{=}0.80$--$0.85$.

  \textbf{Фон против сельхозугодий.} Наиболее частая путаница (схожая текстура
  даже при 448\,px). Устраняется взвешенной кросс-энтропией.

  \textbf{Редкий класс Barren} (4.7\,\% пикселей) — самый низкий IoU из-за
  дисбаланса классов.

  \textbf{Дороги} как тонкие линейные структуры выигрывают от разрешения 448\,px;
  дальнейший рост возможен при переходе на ViT-L/14.
\end{warnbox}

%%=============================================================================
\section{Заключение}

\subsection{Итоги работы}

Разработан и верифицирован полный пайплайн несупервизированной семантической
сегментации снимков ДЗЗ. \textbf{Принципиальное преимущество} предложенного подхода
состоит в том, что он не требует ни одной размеченной выборки на этапе обучения:
семантическая информация извлекается исключительно из структуры данных через
DINOv2-признаки и итеративное самообучение.

\textbf{K-means инициализация} (2\,582\,528 патч-векторов, $d = 768$) создаёт
семантически состоятельную начальную разбивку уже за несколько секунд. Механизм
\textbf{Mean Teacher} с EMA ($\alpha = 0.999$) обеспечивает стабильную генерацию
псевдометок: teacher накапливает исторические знания, student обучается на жёстких
аугментациях. \textbf{Адаптивный порог уверенности} постепенно расширяет множество
учитываемых пикселей по мере роста качества teacher. \textbf{Венгерский алгоритм}
при оценке корректно выравнивает произвольные кластерные ID с семантическими классами
без дополнительного supervision.

\subsection{Ключевые вклады}

\begin{enumerate}
  \item Интеграция DINOv2 + SegFormer для несупервизированной сегментации ДЗЗ.
  \item Трёхфазный пайплайн: K-means $\to$ warm-up $\to$ итеративное self-training.
  \item Верификация на реальном датасете LoveDA: best mIoU \textbf{21.38\,\%}
        (ViT-B/14, 448\,px, RTX~5070~Ti, CUDA~13.2, 2\,ч\,18\,мин).
  \item Публичная кодовая база с воспроизводимыми экспериментами:
        \texttt{make env-gpu \&\& make download \&\& make train-gpu}.
\end{enumerate}

\subsection{Направления дальнейшей работы}

\textbf{Слабая разметка (semi-supervised):} добавление всего 1--5\,\% размеченных
снимков ожидаемо даст скачок mIoU на $+10$--$15\,\%$. \textbf{Мультиспектральные
данные Sentinel-2} (13 каналов): реализованный \texttt{ChannelAdapter} ($13{\to}3$)
готов к тестированию. \textbf{Настройка порога $\theta$} --- переход с $\theta_0{=}0.95$ на
$\theta_0{=}0.80$--$0.85$ позволит активировать consistency loss с первого
раунда и, вероятно, превысит результат warm-up. \textbf{Knowledge distillation:} после обучения teacher
можно дистиллировать в компактную модель для инференса на устройстве без GPU.
\textbf{Temporal consistency:} для серий снимков одного региона временна́я когерентность
служит дополнительным несупервизированным сигналом.

%%=============================================================================
\appendix
\section{Структура проекта}

\begin{table}[H]
\centering
\caption{Модули кодовой базы NNDP-Lab}
\label{tab:structure}
\small
\begin{tabularx}{\textwidth}{L{4.2cm} X}
  \toprule
  \textbf{Файл / Директория} & \textbf{Назначение} \\
  \midrule
  \texttt{configs/default.yaml}   & Все гиперпараметры \\
  \texttt{data/loveda.py}         & Загрузчик LoveDA (train/val, RGB) \\
  \texttt{data/sentinel2.py}      & Загрузчик Sentinel-2 (GeoTIFF, 13 каналов) \\
  \texttt{data/transforms.py}     & Слабые и сильные аугментации \\
  \texttt{models/backbone.py}     & DINOv2 + ChannelAdapter \\
  \texttt{models/decoder.py}      & SegFormer 4-scale MLP decoder \\
  \texttt{models/segmentor.py}    & Segmentor + MeanTeacherModel (EMA) \\
  \texttt{pseudo\_labels/kmeans\_init.py} & Фаза 0: K-means на признаках DINOv2 \\
  \texttt{pseudo\_labels/confidence.py}   & Фильтры уверенности (max-softmax, entropy) \\
  \texttt{training/losses.py}     & MaskedCrossEntropy + ConsistencyLoss (KL) \\
  \texttt{training/trainer.py}    & Полный цикл self-training \\
  \texttt{utils/metrics.py}       & mIoU + Hungarian matching \\
  \texttt{tests/smoke\_test.py}   & Сквозной тест (stub backbone, $<$2 мин) \\
  \texttt{train.py}               & Точка входа обучения \\
  \texttt{evaluate.py}            & Оценка с Hungarian matching \\
  \texttt{Makefile}               & Цели: env, download, train-*, eval, test \\
  \bottomrule
\end{tabularx}
\end{table}

%%=============================================================================
\bibliographystyle{unsrtnat}
\begin{thebibliography}{9}

\bibitem{oquab2024dinov2}
  M.~Oquab \textit{et al.}
  \newblock {DINOv2}: Learning Robust Visual Features without Supervision.
  \newblock \textit{Transactions on Machine Learning Research}, 2024.
  \newblock \href{https://arxiv.org/abs/2304.07193}{arXiv:2304.07193}.

\bibitem{wang2021loveda}
  J.~Wang \textit{et al.}
  \newblock {LoveDA}: A Remote Sensing Land-Cover Dataset for Domain Adaptive Semantic Segmentation.
  \newblock \textit{NeurIPS Datasets and Benchmarks}, 2021.
  \newblock \href{https://arxiv.org/abs/2110.08733}{arXiv:2110.08733}.

\bibitem{tarvainen2017mean}
  A.~Tarvainen and H.~Valpola.
  \newblock Mean teachers are better role models: Weight-averaged consistency targets improve semi-supervised deep learning results.
  \newblock \textit{NeurIPS}, 2017.
  \newblock \href{https://arxiv.org/abs/1703.01780}{arXiv:1703.01780}.

\bibitem{xie2021segformer}
  E.~Xie \textit{et al.}
  \newblock {SegFormer}: Simple and Efficient Design for Semantic Segmentation with Transformers.
  \newblock \textit{NeurIPS}, 2021.
  \newblock \href{https://arxiv.org/abs/2105.15203}{arXiv:2105.15203}.

\bibitem{caron2021dino}
  M.~Caron \textit{et al.}
  \newblock Emerging Properties in Self-Supervised Vision Transformers.
  \newblock \textit{ICCV}, 2021.
  \newblock \href{https://arxiv.org/abs/2104.14294}{arXiv:2104.14294}.

\end{thebibliography}

\end{document}
